% Standalone problem document, generated from the adjacent README.md.
% Compile with XeLaTeX. Mathematical content is maintained in Markdown.
\documentclass[11pt,a4paper]{article}
\usepackage[fontsize=10.5pt]{scrextend}
\usepackage[margin=22mm,headheight=15pt,headsep=9mm,footskip=11mm]{geometry}
\usepackage{fontspec}
\setmainfont{texgyrepagella}[Extension=.otf,UprightFont=*-regular,BoldFont=*-bold,ItalicFont=*-italic,BoldItalicFont=*-bolditalic]
\setsansfont{texgyreheros}[Extension=.otf,UprightFont=*-regular,BoldFont=*-bold,ItalicFont=*-italic,BoldItalicFont=*-bolditalic]
\setmonofont{texgyrecursor}[Extension=.otf,UprightFont=*-regular,BoldFont=*-bold,ItalicFont=*-italic,BoldItalicFont=*-bolditalic,Scale=MatchLowercase]
\usepackage{amsmath,amssymb}
\usepackage{unicode-math}
\setmathfont{texgyrepagella-math.otf}
\usepackage{xcolor,microtype,enumitem,needspace,fancyhdr}
\usepackage{bookmark}
\definecolor{ink}{HTML}{17344B}
\definecolor{accent}{HTML}{197D80}
\definecolor{muted}{HTML}{596773}
\hypersetup{colorlinks=true,linkcolor=accent,urlcolor=accent,pdftitle={IS-01 - Two
permutation matrices generate the doubly stochastic spectral
boundary},pdfauthor={Open Problems in Numerical Linear Algebra}}
\urlstyle{same}
\setlength{\parindent}{0pt}
\setlength{\parskip}{5pt plus 1pt minus 1pt}
\linespread{1.04}
\setlength{\emergencystretch}{3em}
\widowpenalty=10000
\clubpenalty=10000
\displaywidowpenalty=10000
\allowdisplaybreaks[1]
\setlist{nosep,leftmargin=1.5em}
\providecommand{\tightlist}{\setlength{\itemsep}{0pt}\setlength{\parskip}{0pt}}
\providecommand{\pandocbounded}[1]{#1}
\pagestyle{fancy}
\fancyhf{}
\fancyhead[L]{\sffamily\footnotesize\color{muted}OPEN PROBLEMS IN NUMERICAL LINEAR ALGEBRA}
\fancyhead[R]{\sffamily\bfseries\footnotesize\color{accent}IS-01}
\fancyfoot[L]{\parbox[b]{0.9\textwidth}{\sffamily\fontsize{8}{10}\selectfont\color{muted}Editorial ratings; a bounded search cannot certify that a problem remains unsolved.\\Literature check: 2026-09-08}}
\fancyfoot[R]{\sffamily\footnotesize\color{muted}\thepage}
\renewcommand{\headrulewidth}{0.3pt}
\renewcommand{\headrule}{\hbox to\headwidth{\color{accent}\leaders\hrule height \headrulewidth\hfill}}
\makeatletter
\renewcommand{\section}{\Needspace{5\baselineskip}\@startsection{section}{1}{0pt}{12pt plus 2pt minus 2pt}{4pt}{\sffamily\bfseries\large\color{ink}}}
\renewcommand{\subsection}{\Needspace{4\baselineskip}\@startsection{subsection}{2}{0pt}{9pt plus 1pt minus 1pt}{3pt}{\sffamily\bfseries\normalsize\color{ink}}}
\makeatother
\setcounter{secnumdepth}{0}
\begin{document}
{\sffamily\small\color{muted}Eigenvalues and inverse problems\par}
\vspace{3pt}
{\raggedright\hyphenpenalty=10000\exhyphenpenalty=10000\sffamily\bfseries\fontsize{21}{24}\selectfont\color{ink}Two
permutation matrices generate the doubly stochastic spectral
boundary\par}
\vspace{7pt}
{\sffamily\small\textbf{Difficulty:} challenging\quad\textbf{Importance:} interesting
to the community\par}
\vspace{4pt}
{\color{accent}\hrule height 0.8pt}
\vspace{6pt}
\textbf{Status:} open.

\subsection{Problem statement}\label{problem-statement}

Let \(\mathcal D_n\) consist of the real entrywise-nonnegative
\(n\times n\) matrices whose row and column sums are all one, and define
the single-eigenvalue region

\[
D_n=\{z\in\mathbb C:z\in\sigma(A)\text{ for some }A\in\mathcal D_n\}.
\]

Prove or disprove that, for every integer \(n\geq1\) and every
\(z\in\partial D_n\), there exist \(n\times n\) permutation matrices
\(P,Q\) and \(t\in[0,1]\) such that

\[
\det\bigl(zI-tP-(1-t)Q\bigr)=0.
\]

The boundary is taken in the usual topology of \(\mathbb C\), and
\(P=Q\) is allowed. The task concerns one eigenvalue, not simultaneous
realization of a prescribed full spectrum. It would reduce boundary
computation to a finite collection of one-parameter matrix families.
Counterexamples for convex hulls of other matrix groups do not
automatically apply to all permutation matrices.

\subsection{References}\label{references}

Harlev, Johnson, and Lim,
\href{https://arxiv.org/html/1908.03647v2}{\emph{The Doubly Stochastic
Single Eigenvalue Problem: A Computational Approach}}, Conjecture 2.7
and §6; published in \emph{Experimental Mathematics} 31 (2022),
936--945. Verbeken and Ginis,
\href{https://ilas2026.math.vt.edu/docs/ILAS2026-Book-Of-Abstracts.pdf}{ILAS
2026 abstract}, p.~150.

\subsection{Status check}\label{status-check}

Searches for \textit{Harlev Johnson Lim boundary conjecture} and
\textit{doubly stochastic boundary conjecture 2026} found the 2026
authors reporting numerical support through order 25, rather than a
proof. The older Perfect--Mirsky conjecture fails at order five; that
failure does not refute this different assertion.
\end{document}
